\makeabstract
    {
    Water Quality in the Fort River--Mill River Watershed
    }
    {
    Alysha Rodr{\'\i}guez Cab\'an\textsuperscript{1}, Anika Ponnambalam\textsuperscript{1}, Jesse Lebolt\textsuperscript{1}, Anna Martini\textsuperscript{1}
    }
    {
    \textsuperscript{1} Department of Geology, Amherst College, Amherst, MA
    }
    {
    This study aims to identify the spatial and temporal variability of water quality across multiple sites within the Fort River--Mill River watershed in Amherst, Massachusetts. By evaluating chemical and microbial water quality indicators, we can gain insight into the possible environmental stressors affecting the watershed.
    }
    {
    Water samples were collected from various sites in the watershed. Conductivity, alkalinity, and major ions were measured, and the IDEXX Colilert method was used to quantify total coliform and Escherichia coli (E. coli) concentrations.
    }
    {
    Across the four sites sampled for E. coli, concentrations increased following rainfall. At Fearing Brook, E. coli increased from 222.4 MPN/100 mL under dry conditions to 1,986.3 MPN/100 mL after heavy rainfall. Similarly, at the Fort River site, concentrations increased from 151 MPN/100 mL to 488.4 MPN/100 mL following rainfall and later decreased to 275.5 MPN/100 mL after a period of drier conditions. Similar rainfall-related increases were observed at the other two sampling sites. Total coliform concentrations showed positive relationships with chloride and nitrate concentrations. The relationship was moderate for chloride ($R^2 = 0.3266$) and stronger for nitrate ($R^2 = 0.6647$).
    }
    {
    E. coli and total coliform concentrations showed both spatial and temporal variability across the Fort River--Mill River watershed and increased following rainfall events. Total coliform concentrations were also positively associated with chloride and nitrate, suggesting potential links between microbial contamination and nutrient inputs.
    }

    \newpage
    
